\paragraph{From the full cohomology bound to the same determinant scalar.}
Here is the rational divisibility step explicitly. Fix a character
$x$ of $G$, a height-one prime $\mathfrak p$ of $A$, and
$R=A_{\mathfrak p}$. Put $z=F_E(Z_F)_x$, without changing its
normalization. The full global degree-one module is free of rank one
over $R$ by the argument following (K15) and the injectivity of the
maps in (K16). Write it as $Re$, and write $z=\alpha e$.
The higher real cohomology vanishes over $R$, since two is invertible.
Global degree-zero cohomology is zero: an invariant would restrict
to an invariant of the open nonabelian subgroup on the elliptic
representation, on which the scalar coefficient character is trivial.
Thus the positive-cone exact sequence gives
\[
 0\longrightarrow R a\longrightarrow
 H^1(C^+_x\otimes R)\longrightarrow Re\longrightarrow0,
 \qquad
 H^2(C^+_x\otimes R)=H^2(C_{S,x}\otimes R).
\]
The element $a$ here is the image of the fixed real invariant generator
in (E.4), and it is primitive in this positive degree-one module,
because the quotient $Re$ is free. Choose a lift $\widetilde e$ of $e$;
then $(a,\widetilde e)$ is a basis. The equality already proved from
(E.4) becomes
\[
                 u_x^{\rm ev}\theta_x
                  =2\alpha\,a\wedge\widetilde e.
\]
Let $d_1,\ldots,d_r$ be the nonzero Smith entries of a finite free
model of $C^+_x\otimes R$. All its cohomology outside degrees one and
two is zero, its degree-one cohomology is free, and its degree-two
cohomology is torsion. Its inverse determinant lattice, under the
rational identification with $\bigwedge^2H^1$, is therefore generated,
up to an $R$-unit, by
$(\prod_i d_i)a\wedge\widetilde e$. It follows that
\[
 \begin{split}
 v_{\mathfrak p}(u_x^{\rm ev})
 &=v_{\mathfrak p}(\alpha)-\sum_i v_{\mathfrak p}(d_i)\\
 &=\ell_R(H^1(C_{S,x}\otimes R)/Rz)
       -\ell_R H^2(C_{S,x}\otimes R)\ \geq\ 0,
 \end{split}                                                   \tag{E.4a}
\]
where two is a unit and the inequality is exactly (K2).
This uses the full $H^2$ corrected in (K16), rather than only Kato's
$\mathbf H^2_0$.

The intersection of the height-one valuation rings of the normal
domain $A$ gives $u_x^{\rm ev}\in A$ for every $x$. Since $|G|$ is
a power of two, the character idempotents belong to $A[G]$ and these
memberships give $u^{\rm ev}\in A[G]$. This establishes membership
only; it is not a bound for the denominator from character inversion.
The uniform bound comes separately from (E.3): in its free rank-two
$\mathscr O[G]$-module the integral vectors $a,b$ of (E.4) have integral
wedge, so $2^{c_E}u^{\rm ev}\in\mathscr O[G]$ with the fixed exponent
already used for the original lattice. Coefficientwise intersection
in the original group basis gives


\[
              A\cap\mathscr O=\Lambda,\qquad
                  2^{c_E}u^{\rm ev}\in\Lambda[G].                 \tag{E.5}
\]
Indeed an element of $A$ has bounded powers of two as denominators
in its power-series coefficients; membership in $\mathscr O$ makes
their minimum valuation nonnegative. The coefficientwise intersection
retains $c_E$, without a term $v_2|G|$.
\end{proof}
